% Шевът на абстракцията за вход и изход.
%
% Четирите колони са на стъпка 3.25 cm и ширина 3.05 cm, за да се подредят протоколните
% слоеве отгоре и бекендите отдолу точно един под друг. Цялата картина, заедно със
% завъртените етикети отляво, излиза около 13.6 cm и се побира в 14.5 cm без мащабиране.
%
% Черно бяло: трите пласта се различават по фон (бял, сив 25, сив 10), а моделът на
% бекенда се чете по рамката, плътна за завършване и пунктирана за готовност. Никъде
% не се разчита на цвят.
\begin{figure}[htbp]
\centering
\begin{tikzpicture}[
  lane/.style={proc, minimum width=3.05cm, minimum height=0.8cm},
  impl/.style={box, minimum width=3.05cm, minimum height=1.4cm, fill=gray!10, thick},
  compl/.style={impl},                    % модел на завършване
  ready/.style={impl, densely dashed},    % модел на готовност
  seam/.style={box, minimum width=12.8cm, minimum height=1.15cm, fill=gray!25, thick},
  side/.style={rotate=90, font=\footnotesize},
]

% ---- Над шева: това, което не се сменя ------------------------------------
\node[lane, minimum width=12.8cm] (app) at (4.875,5.3)
     {Маршрутизатор и middleware верига};

\node[font=\footnotesize] at (4.875,4.62) {протоколни слоеве};

\node[lane] (h11) at (0.00,3.9) {\code{HTTP/1.1}};
\node[lane] (h2)  at (3.25,3.9) {\code{HTTP/2}};
\node[lane] (h3)  at (6.50,3.9) {\code{HTTP/3}};
\node[lane] (ws)  at (9.75,3.9) {\code{WebSocket}};

\begin{scope}[on background layer]
  \node[grp, fit=(app)(h11)(ws)] (upper) {};
\end{scope}

% ---- Шевът ----------------------------------------------------------------
\node[seam] (seam) at (4.875,2.0)
     {Интерфейс за вход и изход\\[1pt]
      \code{IoContext}\quad\code{Listener}\quad\code{Connection}\quad\code{DatagramSocket}};

\draw[flow] (upper.south) -- (seam.north);

% ---- Под шева: четирите реализации ----------------------------------------
\node[compl] (uring)  at (0.00,0.10) {\code{io\_uring}\\[2pt] Linux\\[1pt] завършване};
\node[ready] (epoll)  at (3.25,0.10) {\code{epoll}\\[2pt] Linux\\[1pt] готовност};
\node[ready] (kqueue) at (6.50,0.10) {\code{kqueue}\\[2pt] macOS, BSD\\[1pt] готовност};
\node[compl] (iocp)   at (9.75,0.10) {\code{IOCP}\\[2pt] Windows\\[1pt] завършване};

% Пунктирана стрелка надолу: реализация на интерфейса, а не извикване през него.
\foreach \n in {uring,epoll,kqueue,iocp}
  \draw[thin flow] (seam.south -| \n.north) -- (\n.north);

% ---- Методологичната бележка ----------------------------------------------
% Скобата обхваща само двата Linux бекенда, защото точно тяхното съжителство
% прави модела независима променлива.
\draw[decorate, decoration={brace, mirror, amplitude=5pt}, semithick]
     ($(uring.south west)+(0,-0.10)$) -- ($(epoll.south east)+(0,-0.10)$);

\node[note, text width=12.3cm, anchor=north] (mnote) at (4.875,-1.25)
     {И двата модела са налични на една и съща система. Моделът на интерфейса,
      готовност или завършване, се мери като независима променлива: една машина,
      един двоичен файл, сменя се само бекендът. На macOS и Windows моделът върви
      заедно с операционната система и двете не могат да се разделят.};

\node[note, text width=12.3cm, anchor=north] at ($(mnote.south)+(0,-0.15)$)
     {Плътна рамка: модел на завършване. Пунктирана рамка: модел на готовност.};

% Завъртените етикети отляво носят цялото послание на фигурата. Центрирани са
% по вертикалната среда на пласта, който именуват, за да не се четат като етикет
% на шева между тях.
\node[side] at (-1.95,4.575) {не се променя};
\node[side] at (-1.95,0.100) {сменя се};

\end{tikzpicture}
\caption{Шев на абстракцията за вход и изход. Над шева стои всичко, което не зависи
от бекенда: маршрутизаторът, middleware веригата и протоколните слоеве. Под шева
стоят четирите реализации. Понеже \code{io\_uring} и \code{epoll} се предлагат на
една и съща операционна система, моделът на интерфейса може да се измери отделно
от нея.}
\label{fig:io_backend_seam}
\end{figure}
